\subsubsection{拉马努金型对称的结构本质：从“坐标对称”到“相位化对称”}\label{subsubsec:ramanujan-phase-lift-symmetry}
本节已将 completion/自对偶诱导的对称性压缩为参数侧的固定几何：unitary slice \(|u|=1\) 及其 Dirichlet 像 \(\Re(s)=\tfrac12\)（注 \ref{rem:xi-critical-line-unit-circle}）。
需要强调的是：此处所谓“对称性”并非若干读出轴之间的坐标置换，而是投影系统自身的自同构（例如 \((z,u)\mapsto(uz,1/u)\) 或 \(u\mapsto 1/u\)）及其固定切片。
在该口径下，\(\Re(s)=\tfrac12\) 的出现并不依赖“对称轴”隐喻，而是相位化闭包在完成化坐标中的唯一自洽切片。
\par
结合圆盘化视界—多极矩接口与 kernel-RH/拉马努金谱界字典（见 \S\ref{subsubsec:selfdual-blaschke-defect-fixed-slice} 末段与附录 \ref{subsubsec:app-horizon-weyl-carath-toeplitz}），可将“拉马努金对称性”作如下结构性重述：它并不在几何上交换坐标轴，而是在协议上强制\textbf{边界酉化/内函数化}，从而以结构方式排除盘内点谱缺陷（即 Blaschke 数据）。
在该意义下，所谓“对称”首先是一条缺陷排斥律：在自对偶（反线性）对称与边界酉化兼容时，指标消失强制 \(B_F\equiv 1\)，点状离切片缺陷只能被推入固定切片的边界承载（定理 \ref{thm:blaschke-defect-zero-fixed-slice-support}）。
同时，圆盘紧化把高虚部的几何贡献压缩到端点 \(w\approx-1\) 的极薄边界层；若存在离临界线零点 \(\rho=\beta+\mathrm{i}\gamma\)（\(\beta<\tfrac12\)、\(\delta=\tfrac12-\beta>0\)），其圆盘像 \(w_\rho=\mathcal{C}(\gamma+\mathrm{i}\delta)\) 满足闭式压缩律
\(
1-|w_\rho|^2=\frac{4\delta}{\gamma^2+(1+\delta)^2}
\)
并因此只能在端点分辨率足够的制度下被核验（定理 \ref{thm:app-jensen-offcritical-kink-localization}）。

\begin{theorem}[相位化提升：\texorpdfstring{$[-2,2]$}{[-2,2]} 谱界 \(\Leftrightarrow\) 纯相位]\label{thm:phase-lift-spectral-bound}
设 \(\mathcal H\) 为复 Hilbert 空间，\(A\) 为 \(\mathcal H\) 上的有界自伴算子。
则下列条件等价：
\begin{enumerate}
  \item \(\mathrm{Spec}(A)\subseteq[-2,2]\)；
  \item 存在酉算子 \(U\) 使得 \(A=U+U^{-1}\)。
\end{enumerate}
并且在 (1) 成立时，可由连续函数演算取
\[
\boxed{\ U:=\exp\!\bigl(\mathrm{i}\,\arccos(A/2)\bigr)\ }
\]
从而 \(A=U+U^{-1}\)。
\end{theorem}
\begin{proof}
\textbf{(2)\(\Rightarrow\)(1).}
若 \(A=U+U^{-1}\) 且 \(U\) 酉，则 \(\mathrm{Spec}(U)\subseteq\TT\)。
由谱映射定理，\(\mathrm{Spec}(A)=f(\mathrm{Spec}(U))\)，其中 \(f(z)=z+z^{-1}\)。
而对 \(z=e^{i\theta}\in\TT\) 有 \(f(z)=2\cos\theta\in[-2,2]\)，故 \(\mathrm{Spec}(A)\subseteq[-2,2]\)。
\par
\textbf{(1)\(\Rightarrow\)(2).}
在 \([-2,2]\) 上定义连续函数 \(g(x):=\exp(\mathrm{i}\arccos(x/2))\)，则 \(|g(x)|=1\)。
令 \(U:=g(A)\) 为由连续函数演算给出的算子，则 \(U\) 为酉算子。
又因 \(g(x)+g(x)^{-1}=2\cos(\arccos(x/2))=x\) 在 \([-2,2]\) 上恒等成立，
故同一恒等式在函数演算下给出
\(
U+U^{-1}=A
\)。
证毕。
\end{proof}
\leanverified{paper\_phase\_lift\_spectral\_bound}

\begin{remark}[Joukowski 压缩的算子论版本]\label{rem:phase-lift-joukowski-operator}
定理 \ref{thm:phase-lift-spectral-bound} 给出命题 \ref{prop:xi-rh-interval-criterion} 的 \((y,S)\)-压缩在算子层面的对应：若 \(Uv=e^{i\theta}v\)，则 \((U+U^{-1})v=2\cos\theta\,v\)。
因此 \(\mathrm{Spec}(U)\subseteq\TT\) 经 Joukowski 映射压缩到 \([-2,2]\)；反向则把 \([-2,2]\) 上的谱提升为单位圆相位。
这一“压缩 \(\Leftrightarrow\) 提升”的同型性，在本节的完成化变量 \(S=y+y^{-1}\) 与附录的半角完成化变量 \(s=\mathfrak{r}+\mathfrak{r}^{-1}\) 中同步出现（注 \ref{rem:xi-joukowski-toy-rh-alignment} 与注 \ref{rem:xi-half-angle-endpoint-paradigm}）。
\end{remark}

\begin{definition}[幅度缺陷：极分解中的幅度负担]\label{def:amp-defect}
设 \(T\) 为复 Hilbert 空间 \(\mathcal H\) 上的可逆有界算子，作极分解
\[
T=V|T|,\qquad |T|:=(T^\ast T)^{1/2}\succ 0,
\]
其中 \(V\) 为酉算子。
定义 \(T\) 的\textbf{幅度缺陷}为
\[
\boxed{\ \mathfrak{d}_{\mathrm{amp}}(T):=\bigl\|\log|T|\bigr\|_{\mathrm{op}}\ \in[0,\infty)\ },
\]
其中 \(\log|T|\) 由连续函数演算定义。
\end{definition}

\begin{proposition}[纯相位演化 \(\Leftrightarrow\) 幅度缺陷为零]\label{prop:amp-defect-unitary}
\leanverified{paper\_amp\_defect\_unitary}
在定义 \ref{def:amp-defect} 的记号下，有
\[
\boxed{\ \mathfrak{d}_{\mathrm{amp}}(T)=0\quad\Longleftrightarrow\quad |T|=I\quad\Longleftrightarrow\quad T\ \text{为酉算子}\ }.
\]
\end{proposition}
\begin{proof}
\(\mathfrak{d}_{\mathrm{amp}}(T)=0\) 等价于 \(\log|T|=0\)，从而 \(|T|=I\)。
此时 \(T=V|T|=V\) 且 \(V\) 为酉算子，故 \(T\) 酉。
反向蕴含显然：若 \(T\) 酉，则 \(|T|=(T^\ast T)^{1/2}=I\)，从而 \(\mathfrak{d}_{\mathrm{amp}}(T)=0\)。
证毕。
\end{proof}

\begin{remark}[与三轴地址语义的对齐]\label{rem:amp-defect-address-axis}
在本文“地址先于取值”的三轴地址打包 \(a=(a_{\mathrm{vis}},a_{\mathrm{res}},a_{\mathrm{mode}})\) 中（定义 \ref{def:multi_axis_address}），径向/幅度通道属于模式轴 \(a_{\mathrm{mode}}\)。
因此，若读出制度将可见域固定为相位闭包并拒绝外置该模式轴，则任意 \(\mathfrak{d}_{\mathrm{amp}}(T)>0\) 的演化都不可能在闭环内维持可逆核验：幅度信息必须作为正交证书外置，否则只能退化为 \(\mathsf{NULL}\)（见定理 \ref{thm:xi-orthogonal-slice-structure} 的协议正交条目）。
\end{remark}

\begin{theorem}[拉马努金型相位化原则（制度内充分条件）]\label{thm:ramanujan-phase-lift-principle}
设 \(\{\widetilde T_\alpha\}\) 为一族在复 Hilbert 空间 \(\mathcal H\) 上实现的有界自伴算子，并满足谱界
\[
\mathrm{Spec}(\widetilde T_\alpha)\subseteq[-2,2]\qquad(\forall \alpha).
\]
则对每个 \(\alpha\) 存在酉算子 \(U_\alpha\) 使得
\[
\boxed{\ \widetilde T_\alpha=U_\alpha+U_\alpha^{-1}\ }.
\]
因此，上述谱界把可能携带幅度信息的自由度压缩为纯相位：在拒绝引入连续幅度寄存器的纯相位核验制度中，\(\widetilde T_\alpha\) 可被视为相位算子族 \(\{U_\alpha\}\) 的可见像（命题 \ref{prop:amp-defect-unitary}）。
\end{theorem}
\begin{proof}
对每个 \(\alpha\) 直接应用定理 \ref{thm:phase-lift-spectral-bound} 即得 \(U_\alpha\) 的存在与显式构造；而 \(U_\alpha\) 酉蕴含其幅度缺陷为零（命题 \ref{prop:amp-defect-unitary}）。证毕。
\end{proof}

\begin{remark}[与核版 RH/GRH 的同型性：半径界 \texorpdfstring{$\Leftrightarrow$}{<->} 相位化]\label{rem:ramanujan-phase-lift-kernel}
在有限维核口径中，“Ramanujan/RH-sharp 半径界”以 \(|\mu|\le\sqrt{\lambda}\) 的形式出现（例如定理 \ref{thm:pom-rh-sharp}；并见注 \ref{rem:pom-profinite-ihara-ramanujan} 的 Ihara--Ramanujan 同型解释）。
将其写成归一化相位 \(y=\mu/\sqrt{\lambda}\) 后，\(|y|=1\) 的达界模态经 Joukowski 压缩 \(S=y+y^{-1}\) 落入 \([-2,2]\)（注 \ref{rem:xi-joukowski-toy-rh-alignment}），从而在完成化坐标中体现为“可相位化”的结构性约束。
因此，拉马努金语境中的对称性并不对应“坐标轴可互换”，而对应于对称算子族被强制推入相位化闭包所诱导的纯相位可核验性。
\end{remark}

\begin{proposition}[酉行列式的零点定位：\texorpdfstring{$|u|=1$}{|u|=1}]\label{prop:unitary-determinant-zero-unit-circle}
设 \(U\in\CC^{n\times n}\) 为酉矩阵，并令
\[
D(u):=\det(I-uU)\qquad(u\in\CC).
\]
若 \(D(u)=0\)，则 \(|u|=1\)。
进一步，对任意 \(L>1\) 令 \(u_L(s)=L^{2s-1}\)，则
\[
D(u_L(s))=0\quad\Longrightarrow\quad \Re(s)=\tfrac12.
\]
\end{proposition}
\begin{proof}
\(D(u)=0\) 当且仅当 \(\det(I-uU)=0\)，等价于 \(1/u\in\mathrm{Spec}(U)\)。
而酉矩阵的全部谱点满足模长为 \(1\)，故 \(|1/u|=1\)，即 \(|u|=1\)。
对 \(u=u_L(s)=L^{2s-1}\)，有 \(|u|=L^{2\Re(s)-1}\)，故 \(|u|=1\) 等价于 \(\Re(s)=\tfrac12\)。
证毕。
\end{proof}
\leanverified{paper\_unitary\_determinant\_zero\_unit\_circle}

\begin{theorem}[相位化闭包的固定切片性：可核验零点的制度性约束]\label{thm:phase-lift-fixed-slice}
在 unitary-slice 读出制度中，若某个核验对象可在有限维可审计口径下表示为
\[
\mathcal{D}(s):=\det\!\bigl(I-u_L(s)U\bigr),
\qquad
u_L(s)=L^{2s-1},\ L>1,
\]
其中 \(U\) 为酉矩阵，且读出制度不在模式轴外置径向/幅度寄存器，
则 \(\mathcal{D}\) 的全部零点事件都满足 \(\Re(s)=\tfrac12\)。
等价地，\(\Re(s)=\tfrac12\) 是相位化闭包在完成化坐标中的唯一固定切片。
\end{theorem}
\begin{proof}
由命题 \ref{prop:unitary-determinant-zero-unit-circle}，\(\mathcal{D}(s)=0\Rightarrow |u_L(s)|=1\Rightarrow \Re(s)=\tfrac12\)。
而“拒绝外置径向/幅度寄存器”的制度假设确保该零点事件无需额外模式轴证书即可在相位闭包内被核验，从而上述定位即为制度内的零点切片约束。证毕。
\end{proof}
\leanverified{paper\_phase\_lift\_fixed\_slice}

\begin{proposition}[有限型 Hilbert--P{\'o}lya：\texorpdfstring{$\Phi_L$}{PhiL} 的酉谱实现与 CMV 规范型]\label{prop:xi-hilbert-polya-cmv}
在推论 \ref{cor:xi-rh-unit-circle-reciprocal} 的记号下，令 \(d:=\deg P_L\) 并记
\[
\Phi_L(y)=y^dP_L(y+y^{-1})\in\RR[y],
\qquad
D:=\deg \Phi_L=2d.
\]
令 \(\widetilde\Phi_L(y):=\Phi_L(y)/\mathrm{lc}(\Phi_L)\) 为首一化（monic）版本（不改变零点集合）。
则：
\begin{enumerate}
  \item 若 \(\widetilde\Phi_L\) 的全部根为 \(\{e^{i\theta_j}\}_{j=1}^{D}\subset\TT\)，则存在显式酉矩阵
  \[
  U_L^{\mathrm{diag}}:=\mathrm{diag}(e^{i\theta_1},\dots,e^{i\theta_D})
  \]
  使得
  \[
  \boxed{\ \widetilde\Phi_L(y)=\det(yI-U_L^{\mathrm{diag}})\ }.
  \]
  因而 \(\Xi\)-RH（等价于 \(\Phi_L\) 的根全落在 \(|y|=1\)）可被表述为一个\textbf{酉算子谱问题}：零点相位即为酉谱。
  \item 更进一步，存在一个五对角 CMV-like 伴随矩阵 \(U_L^{\mathrm{CMV}}\in\CC^{D\times D}\) 使得
  \[
  \boxed{\ \widetilde\Phi_L(y)=\det(yI-U_L^{\mathrm{CMV}})\ }.
  \]
  且该伴随矩阵满足如下“酉性 \(\Leftrightarrow\) 单位圆根”判别：\(U_L^{\mathrm{CMV}}\) 为酉当且仅当 \(\widetilde\Phi_L\) 的全部根满足 \(|y|=1\)。
  在此构造中，\(U_L^{\mathrm{CMV}}\) 可由一列 Verblunsky/Schur 参数 \(\alpha_0(L),\dots,\alpha_{D-1}(L)\) 表达；这些参数可由 \(\widetilde\Phi_L\) 的系数经 Schur 递归（反射系数算法）显式得到，并满足酉性的等价约束
  \[
  \boxed{\ |\alpha_k(L)|<1\ (0\le k\le D-2),\qquad |\alpha_{D-1}(L)|=1\ }.
  \]
\end{enumerate}
\end{proposition}
\begin{proof}
第一条为谱分解的直接构造：取 \(\widetilde\Phi_L\) 的根写作 \(e^{i\theta_j}\)，则 \(\det(yI-U_L^{\mathrm{diag}})=\prod_{j=1}^{D}(y-e^{i\theta_j})=\widetilde\Phi_L(y)\)。
第二条为 CMV-like 伴随矩阵结论。任意首一多项式都可被实现为某个 CMV-like Hessenberg/五对角伴随矩阵的特征多项式，并且该伴随矩阵可用 Schur（Verblunsky）参数显式参数化；在该参数化下，酉性等价于参数落在单位圆盘/单位圆的约束，从而等价于特征多项式根分布在单位圆上（参见 \cite{CanteroMoralVelazquez2003CMV,Simon2005OPUC1,BevilacquaDelCorsoGemignani2015CMVCompanion}）。证毕。
\end{proof}

\begin{remark}[Cayley--Hilbert--P{\'o}lya 合并：酉谱与自伴谱的同一对象]\label{rem:xi-cayley-hilbert-polya-unify}
命题 \ref{prop:xi-rh-cayley-hyperbolicity} 把“单位圆根条件”完全转译为“实轴全实根条件”。
与之配套，若 \(U\) 为酉矩阵且 \(1\notin\mathrm{Spec}(U)\)，则 Cayley 变换
\[
\boxed{\ H:=\mathrm{i}(I+U)(I-U)^{-1}\ }
\]
给出自伴矩阵（且 \(\mathrm{Spec}(H)\subset\RR\)）。
更具体地，若 \(Uv=yv\) 且 \(|y|=1\)（\(y\neq 1\)），则同一特征向量满足
\[
Hv=xv,\qquad x=\mathrm{i}\,\frac{1+y}{1-y}=\mathcal{C}^{-1}(y)\in\RR,
\]
这与 \eqref{eq:xi-cayley-y-from-x} 的根坐标变换一致。
因此在 Cayley 坐标下，命题 \ref{prop:xi-hilbert-polya-cmv} 的酉谱实现可等价地理解为一套自伴谱实现：酉 Hilbert--P{\'o}lya 与自伴 Hilbert--P{\'o}lya 在此并非并行接口，而是同一对象的两种坐标读出（\(y\in\TT\) 与 \(x\in\overline{\RR}\)）。
当 \(\widetilde\Phi_L\) 在 \(y=1\) 处出现零点时，对应 \(x=\infty\) 的射影根，等价地表现为 \(I-U\) 的退化；此时可用射影紧化语言把上述对应关系无歧义延拓。
\end{remark}

\begin{remark}[逆谱接口：从临界线相位到 CMV/Schur 参数]\label{rem:xi-cmv-inverse-spectral}
命题 \ref{prop:xi-hilbert-polya-cmv} 的意义不仅在于把“单位圆根”解释为“酉谱”，还在于提供一组可计算的\textbf{参数化坐标系}：
CMV-like 伴随矩阵 \(U_L^{\mathrm{CMV}}\) 的自由度可由 Verblunsky/Schur 参数 \((\alpha_k)\) 递归编码，且酉性等价于 \(|\alpha_k|<1\)（末项落在单位圆）这一组显式约束。
因此，将 \(\Xi_{\Delta,L}\) 的临界线零点结构规定为某种目标相位集合（或相位分布）的任务，可转写为一个有限维\textbf{酉矩阵逆谱/参数拟合问题}：
先在 \((\alpha_k)\) 坐标中构造（或拟合）满足酉性约束的参数序列，再由 Schur 递归恢复 \(\widetilde\Phi_L\) 及其相位零点。
该接口与附录 \ref{app:real-input-40-finite-rh} 中 toy-explicit-formula 的有限窗零点设计任务在结构上同型。
\end{remark}

\begin{theorem}[有限 CMV 证书的点态稳定化 \texorpdfstring{$\Rightarrow$}{=>} 无限维 CMV 谱实现]\label{thm:xi-cmv-stabilization-infinite}
设 \(\{U_L^{\mathrm{CMV}}\}_{L}\) 为一族有限维 CMV-like 伴随矩阵实现（命题 \ref{prop:xi-hilbert-polya-cmv}），其维数记为 \(D_L\to\infty\)。
记相应的 Verblunsky/Schur 参数为
\[
\alpha_0(L),\dots,\alpha_{D_L-1}(L),
\qquad
|\alpha_k(L)|<1\ (0\le k\le D_L-2),\ |\alpha_{D_L-1}(L)|=1.
\]
若对每个固定指标 \(n\ge 0\)，存在 \(L_0(n)\) 使得当 \(L\ge L_0(n)\) 时 \(\alpha_n(L)\) 与 \(L\) 无关，并记其稳定值为 \(\alpha_n\)，则：
\begin{enumerate}
  \item 存在唯一的无限维 CMV 幺正算子 \(U_\infty\) 作用在 \(\ell^2(\ZZ_{\ge 0})\) 上，其 Verblunsky 系数序列为 \((\alpha_n)_{n\ge 0}\)（特别地，对每个 \(n\) 有 \(|\alpha_n|<1\)）。
  \item 若把 \(U_L^{\mathrm{CMV}}\) 视为 \(\ell^2(\ZZ_{\ge 0})\) 上的幺正算子
  \[
  \widetilde U_L:=U_L^{\mathrm{CMV}}\oplus I_{\ell^2(\{D_L,D_L+1,\dots\})},
  \]
  则对任意有限支撑向量 \(\xi\) 有 \(\widetilde U_L\xi\to U_\infty\xi\)；换言之，\(\widetilde U_L\to U_\infty\) 在强算子拓扑下收敛（并且在有限支撑向量上最终恒等）。
  \item 令 \(\mu_L,\mu_\infty\) 分别为 \(\widetilde U_L\) 与 \(U_\infty\) 关于基向量 \(e_0\) 的谱测度，则 \(\mu_L\Rightarrow \mu_\infty\)（弱收敛）。
  等价地，Carath\'eodory 函数
  \[
  F_L(z):=\int_{\TT}\frac{w+z}{w-z}\,d\mu_L(w)\qquad(|z|<1)
  \]
  在单位圆盘内按紧子集局部一致收敛到 \(F_\infty\)。
\end{enumerate}
\end{theorem}

\begin{proof}
关于 (1)：对任意满足 \(|\alpha_n|<1\) 的序列 \((\alpha_n)_{n\ge 0}\)，存在唯一与之对应的无限 CMV 幺正算子（见 \cite{Simon2005OPUC1}）。

关于 (2)：CMV 矩阵为五对角形式，其前若干行/列的矩阵元仅依赖于有限个 Verblunsky 系数 \(\alpha_0,\dots,\alpha_{N}\)。
给定有限支撑向量 \(\xi\)，取 \(N\) 使得 \(\mathrm{supp}(\xi)\subset\{0,1,\dots,N\}\)。由点态稳定化，存在 \(L\) 足够大使得 \(D_L>N+2\) 且 \(\alpha_k(L)=\alpha_k\) 对 \(0\le k\le N+1\) 成立；于是 \(\widetilde U_L\) 与 \(U_\infty\) 在前 \(N+2\) 个基向量张成的子空间上的矩阵元一致，从而 \(\widetilde U_L\xi=U_\infty\xi\)。因此得到强收敛。

关于 (3)：由 (2) 得到对每个整数 \(m\ge 0\)，有矩 \(\langle e_0,\widetilde U_L^{\,m}e_0\rangle\to \langle e_0,U_\infty^{\,m}e_0\rangle\)。
这等价于 \(\mu_L\) 的全部三角多项式矩按 \(L\to\infty\) 收敛到 \(\mu_\infty\) 的矩；由三角多项式在 \(C(\TT)\) 中稠密即得谱测度弱收敛。
Carath\'eodory 函数在 \(|z|<1\) 内具有幂级数展开，其系数由上述矩决定，故可推出 \(F_L\) 的局部一致收敛。证毕。
\end{proof}


\begin{conclusion}[有限窗谱矩的终端冻结：点态稳定化把无穷维读出降为有限证书]\label{con:xi-cmv-finite-window-freeze}
在定理 \ref{thm:xi-cmv-stabilization-infinite} 的设定下，
对任意 \(N\ge 0\) 与任意整数 \(m\ge 0\)，存在 \(L_\ast(N,m)\) 使对所有 \(L\ge L_\ast(N,m)\) 有
\[
\langle e_i,\widetilde U_L^{\,m}e_j\rangle=\langle e_i,U_\infty^{\,m}e_j\rangle
\qquad(0\le i,j\le N).
\]
因此任意有限窗口上的有限阶谱矩在某一有限级别的 CMV 证书处即已完全冻结。
等价地，Carath\'eodory 函数 \(F_L\) 在 \(0\) 处的任意有限阶 Taylor 系数在某一有限级别处稳定并可离线核验。
\end{conclusion}
\begin{proof}
CMV 矩阵为五对角带状算子，故对每个固定 \(j\) 与 \(m\)，向量 \(U_\infty^{\,m}e_j\) 具有有限支撑；
取 \(K\) 使得对所有 \(0\le j\le N\) 均有 \(\mathrm{supp}(U_\infty^{\,m}e_j)\subseteq\{0,1,\dots,K\}\)。
由定理 \ref{thm:xi-cmv-stabilization-infinite} 的点态稳定化，
对每个 \(k\le K+1\) 可取 \(L\) 足够大使 \(\alpha_k(L)=\alpha_k\) 且 \(D_L>K+2\)，
从而 \(\widetilde U_L\) 与 \(U_\infty\) 在 \(\mathrm{span}\{e_0,\dots,e_K\}\) 上的矩阵元一致。
因此对 \(0\le j\le N\) 有 \(\widetilde U_L^{\,m}e_j=U_\infty^{\,m}e_j\)，再取内积即得所示谱矩冻结。
最后，\(F_L\) 的幂级数系数由谱测度矩确定（定理 \ref{thm:xi-cmv-stabilization-infinite} 的第 (3) 条），故其有限阶 Taylor 系数亦在同一有限级别处稳定。证毕。
\end{proof}

\endinput

\begin{remark}[证书化验零：从 \texorpdfstring{$P_L$}{PL} 或 \texorpdfstring{$\Phi_L$}{PhiL} 到 Hardy 截面零点区间]\label{rem:xi-hard-zero-certificate}
在完成化表示下，
\[
\Xi_{\Delta,L}\!\left(\tfrac12+it\right)=P_L\!\bigl(2\cos(t\log L)\bigr),
\]
因此临界线零点问题可完全压缩为一元多项式根的证书化问题。
例如可采用如下“可机核验”的有限流程：先用 Sturm 序列（或等价的有理区间算术）把 \(P_L\) 在 \([-2,2]\) 内的实根隔离成互不相交的有理区间；再把每个根区间通过单调的 \(\arccos(\cdot/2)\) 映射为一个 \(t\)-零点区间并按周期 \(2\pi/\log L\) 平移复制；最后用三角多项式的显式导数界在每个区间内给出“至多一个零点”的唯一性证书。
同步核加权例子中，脚本 \texttt{scripts/exp\_sync\_kernel\_xi\_cubic\_cos\_polynomial.py} 已直接输出 \(P_L\) 的根与对应离线竖线参数，可作为该证书流程的最小可审计实例。
\end{remark}

\begin{proposition}[Sturm--Jacobi 规范化：\texorpdfstring{$P_L$}{PL} 的自伴三对角实现]\label{prop:xi-hilbert-polya-jacobi}
沿用命题 \ref{prop:xi-rh-interval-criterion} 的记号，令 \(P_L(S)\in\RR[S]\) 并记其首一化版本
\[
\widetilde P_L(S):=\frac{P_L(S)}{\mathrm{lc}(P_L)}\qquad(\text{首项系数归一化为 }1).
\]
若 \(\widetilde P_L\) 的全部根均为实数（特别地，若 \(\Xi\)-RH 成立则必然如此），则存在一个实对称三对角矩阵
\[
J_L=\begin{pmatrix}
b_1 & a_1 \\
a_1 & b_2 & a_2\\
& a_2 & \ddots & \ddots\\
& & \ddots & b_{d-1} & a_{d-1}\\
& & & a_{d-1} & b_d
\end{pmatrix},\qquad a_k>0,\ b_k\in\RR,\ d:=\deg P_L,
\]
使得
\[
\boxed{\ \widetilde P_L(S)=\det(SI-J_L)\ }.
\]
并且“根区间判据”可改写为自伴谱条件
\[
\boxed{\ \mathrm{Roots}(P_L)\subseteq[-2,2]\quad\Longleftrightarrow\quad \mathrm{Spec}(J_L)\subseteq[-2,2]\quad\Longleftrightarrow\quad \|J_L\|_{\mathrm{op}}\le 2\ }.
\]
此外，\(J_L\) 可由 \(P_L\) 的 Euclid--Sturm 链（或等价的 Stieltjes 连分式数据）构造得到：换言之，Sturm 序列不仅给出“根落在区间”的判定，也同时给出一个\textbf{规范化的有限维 Hilbert--P{\'o}lya 自伴算子候选}。
\end{proposition}
\begin{proof}
Jacobi 三对角实现及其由 Euclid--Sturm 链/连分式数据的构造见 \cite{Fiedler1990SymmetricCharpoly,Schmeisser1993TridiagonalGivenCharpoly}；等价性由 \(\mathrm{Spec}(J_L)=\mathrm{Roots}(\det(SI-J_L))\) 与自伴谱界 \(\mathrm{Spec}(J_L)\subset[-2,2]\iff \|J_L\|_{\mathrm{op}}\le 2\) 直接推出。证毕。
\end{proof}

\begin{remark}[Artin 通道化的完成化判据：把核版 (G)RH 拆成有限个单位圆/区间根检验]\label{rem:xi-artin-channel-pipeline}
第 \ref{sec:A-kernel-compare-main} 节的定义 \ref{def:kernel-group-twist} 引入了带有限群标签的通道矩阵族 \(M_{K,\rho}(u)\) 与通道因子
\(
L_{K,\rho}(z,u):=\det(I-zM_{K,\rho}(u))^{-1}
\)，并在命题 \ref{prop:kernel-artin-factorization} 给出覆盖行列式的 Artin 因子化。
若某个核的自对偶/完成化结构与群作用相容，使得对每个不可约表示 \(\rho\) 的通道行列式
\[
\Delta_\rho(z,u):=\det(I-zM_{K,\rho}(u))
\]
仍满足二变量自对偶 \(\Delta_\rho(z,u)=\Delta_\rho(uz,1/u)\)，并可写成完成化形式
\[
\Delta_\rho(z,u)=\widehat\Delta_\rho\!\left(z\sqrt{u},\ \sqrt{u}+\frac{1}{\sqrt{u}}\right),
\]
则对任意底 \(L>1\) 可定义通道截面
\[
\Xi_{\rho,L}(s):=\Delta_\rho(L^{-s},L^{2s-1})
=\widehat\Delta_\rho\!\bigl(L^{-1/2},S_L(s)\bigr),
\qquad S_L(s)=L^{s-1/2}+L^{1/2-s}.
\]
从而命题 \ref{prop:xi-rh-interval-criterion} 与推论 \ref{cor:xi-rh-unit-circle-reciprocal} 逐通道给出等价判据：
\[
\Xi_{\rho,L}\ \text{的全部零点在}\ \Re(s)=\tfrac12
\quad\Longleftrightarrow\quad
P_{L,\rho}(S):=\widehat\Delta_\rho(L^{-1/2},S)\ \text{的全部根在}\ [-2,2].
\]
因此，“完成化 + Artin 因子化”把核版 (G)RH 的检查压缩为有限多个一元多项式根的证书化问题（Sturm 序列/区间算术即可独立核验）；并可与第 \ref{sec:A-kernel-compare-main} 节的通道级谱半径界（定义 \ref{def:kernel-rh}）形成一条可执行的检验流水线。
\end{remark}

\paragraph{条带版 \texorpdfstring{$\Xi$}{Xi}-RH：只约束开临界条带内的离线零点}
RH 的语义并非“全体零点皆在临界线”，而是只约束开临界条带 \(0<\Re(s)<1\) 内的非平凡零点；
条带外的零点被视作平凡零点。
对 \(\Xi_{\Delta,L}\) 也同样有一个自然的“条带版”核判据：在变量 \(y=L^{s-1/2}\) 下，
开临界条带恰对应于环域 \(L^{-1/2}<|y|<L^{1/2}\)，从而边界 \(\Re(s)=0,1\) 在迹坐标 \(S=y+y^{-1}\) 上塌缩为一个显式禁区窗口。

\begin{proposition}[条带内离线零点的禁区判据（实根情形）]\label{prop:xi-strip-forbidden-real}
在命题 \ref{prop:xi-rh-interval-criterion} 的完成化表示下，仍令
\[
P_L(S):=\widehat\Delta\!\bigl(L^{-1/2},S\bigr)\in\RR[S],
\]
并额外假设 \(P_L\) 的全部根均为实数。
则下述两条陈述等价：
\begin{enumerate}
  \item 存在 \(\Xi_{\Delta,L}\) 的零点 \(s\) 满足 \(0<\Re(s)<1\) 且 \(\Re(s)\neq \tfrac12\)；
  \item \(P_L\) 在禁区
  \[
  (2,\sqrt{L}+\tfrac{1}{\sqrt{L}})\ \cup\ (-\sqrt{L}-\tfrac{1}{\sqrt{L}},-2)
  \]
  上有根。
\end{enumerate}
\end{proposition}
\begin{proof}
令 \(y:=L^{s-1/2}\)。则 \(\Xi_{\Delta,L}(s)=0\) 当且仅当 \(S_0:=y+y^{-1}\) 是 \(P_L\) 的某个根（见命题 \ref{prop:xi-rh-interval-criterion} 的证明）。
当 \(S_0\in\RR\setminus[-2,2]\) 时，二次方程 \(y+y^{-1}=S_0\) 的两根为
\[
y_\pm=\frac{S_0\pm\sqrt{S_0^2-4}}{2},\qquad y_+y_-=1,
\]
且满足 \(|y_+|\ge 1\ge |y_-|\)。
又
\[
0<\Re(s)<1\quad\Longleftrightarrow\quad L^{-1/2}<|y|<L^{1/2},
\]
因此出现条带内离线零点当且仅当对某个根 \(S_0\) 有 \(1<|y_+|<\sqrt{L}\)。
对实根情形，映射 \(y\mapsto y+y^{-1}\) 在区间 \((1,\sqrt{L})\) 上严格递增，并把端点映到 \(2\) 与 \(\sqrt{L}+1/\sqrt{L}\)；
同理在 \((-\sqrt{L},-1)\) 上严格递增并把端点映到 \(-\sqrt{L}-1/\sqrt{L}\) 与 \(-2\)。
于是等价性成立。
\end{proof}

\begin{remark}[复根情形：禁区是椭圆域而非仅两段实区间]\label{rem:xi-strip-ellipse-forbidden}
若 \(P_L\) 允许复根，则“条带内离线零点”对应于根落入 Joukowski 变换 \(S=y+y^{-1}\) 把圆 \(|y|=\sqrt{L}\) 映出的椭圆内部（并避开 \([-2,2]\)）：
\[
S(\theta)=\Bigl(\sqrt{L}+\frac{1}{\sqrt{L}}\Bigr)\cos\theta
+i\Bigl(\sqrt{L}-\frac{1}{\sqrt{L}}\Bigr)\sin\theta,\qquad \theta\in[0,2\pi).
\]
上面的两段实禁区只是该椭圆与实轴的交。
\end{remark}

\begin{remark}[同步核加权情形的三次余弦多项式与离线零点列]\label{rem:sync-kernel-xi-cubic}
对同步核加权情形有 \(\lambda=3\)，并且完成化行列式 \(\widehat\Delta(w,S)\) 具有整系数闭式（命题 \ref{prop:sync-kernel-weighted-hat-delta}）。
沿自对偶截面 \(z=3^{-s},u=3^{2s-1}\) 得 \(w=3^{-1/2}\)，于是临界线上的 Hardy 截面塌缩为一个三次余弦多项式（可复算输出）：
% Auto-generated; do not edit by hand.
\begin{equation}\label{eq:sync-kernel-xi-cubic-cos}
\boxed{
27\,\Xi_{\Delta,3}\!\left(\tfrac12+it\right)=\sqrt{3} S^{3} - 2 S^{2} - 6 \sqrt{3} S - 4,\qquad S:=2\cos(t\log 3).
}
\end{equation}

其三根见表 \ref{tab:sync_kernel_xi_cubic_roots}：
\begin{table}[H]
\centering
\scriptsize
\setlength{\tabcolsep}{6pt}
\caption{Roots of the cubic completion slice polynomial $P(S)=\widehat\Delta(3^{-1/2},S)$ (reported from $27\,P(S)$; scaling does not affect root locations).}
\label{tab:sync_kernel_xi_cubic_roots}
\begin{tabular}{r r r}
\toprule
$i$ & $S_i$ & $S_i\in[-2,2]$?\\
\midrule
1 & -1.642380022755384 & yes\\
2 & -0.435047917672508 & yes\\
3 & 3.232128478807143 & no\\
\bottomrule
\end{tabular}
\end{table}

由于存在根 \(S_0>2\)，由命题 \ref{prop:xi-rh-interval-criterion} 立刻得到：全局意义下该 \(\Xi_{\Delta,3}\) 的零点并非全落在临界线（存在离线竖线）。
然而该核的三根均为实数，且最大的根满足 \(S_0>\sqrt{3}+\tfrac{1}{\sqrt{3}}\)，因此由命题 \ref{prop:xi-strip-forbidden-real} 可知：开临界条带 \(0<\Re(s)<1\) 内不会出现离线零点。
对应的两条离线竖线可写成
\[
\Re(s)=\tfrac12\pm\sigma_0,\qquad
\sigma_0:=\frac{\log\bigl((S_0+\sqrt{S_0^2-4})/2\bigr)}{\log 3},
\]
其数值（以及虚部步长 \(\Delta_t=2\pi/\log 3\)）见可复现输出：
% Auto-generated; do not edit by hand.
\begin{equation}\label{eq:sync-kernel-xi-offcritical-shift}
\boxed{
S_0\approx 3.232128478807143,\qquad\sigma_0\approx 0.964603276500941,\qquad\Delta_t:=\frac{2\pi}{\log 3}\approx 5.719201734760254.
}
\end{equation}

并且虚部步长恒为 \(\Delta_t=2\pi/\log 3\)（由 \(y=L^{s-1/2}\) 的对数多值性给出）。
\end{remark}

\paragraph{单尺度 Dirichlet 化的虚周期刚性：从“梳齿晶格”到零点密度障碍}
在有限维/单尺度口径下，把复变量 \(s\) 压到单一指数 \(z=L^{-s}\)（或其同一底的有限个指数幂）会强制出现严格虚周期，从而把零点/极点限制在有限条等差竖线簇上。

\begin{proposition}[单尺度 Dirichlet 化的严格虚周期]\label{prop:single-scale-dirichlet-imag-period}
\leanverified{paper\_single\_scale\_dirichlet\_imag\_period}
取常数 \(L>1\)，并令 \(g\) 为一变量亚纯函数。定义
\[
f(s):=g(L^{-s}).
\]
则 \(f\) 在虚方向具有严格周期
\[
T_L:=\frac{2\pi i}{\log L},\qquad f(s+T_L)=f(s).
\]
特别地，对任意有限维矩阵 \(M\) 定义
\(
\widehat{\zeta}_{M,L}(s):=\det(I-L^{-s}M)^{-1}
\)，则有 \(\widehat{\zeta}_{M,L}(s+T_L)=\widehat{\zeta}_{M,L}(s)\)。
\end{proposition}
\begin{proof}
由 \(L^{-(s+T_L)}=\exp(-(s+T_L)\log L)=\exp(-s\log L)\exp(-2\pi i)=L^{-s}\) 直接得出。
\end{proof}

\begin{corollary}[零/极点的梳齿晶格与线性计数上界]\label{cor:single-scale-dirichlet-lattice}
\leanverified{paper\_single\_scale\_dirichlet\_lattice}
在命题 \ref{prop:single-scale-dirichlet-imag-period} 的记号下，若 \(g\) 为有理函数，则：
\begin{enumerate}
  \item \(f\) 的全部零与极点在虚方向上按步长 \(2\pi/\log L\) 周期复制；更具体地，它们是有限多个等差竖线簇的并。
  \item 在任意固定竖直条带 \(\sigma_1\le \Re(s)\le \sigma_2\) 内，满足 \(|\Im(s)|\le T\) 的零点与极点总数为 \(O(T)\)。
\end{enumerate}
\end{corollary}
\begin{proof}
由于 \(g\) 的零点/极点在 \(z\)-平面是有限集 \(\{z_j\}\)（按重数计），解方程 \(L^{-s}=z_j\) 得
\[
s=\frac{-\log z_j+2\pi i k}{\log L},\qquad k\in\ZZ,
\]
从而每个 \(z_j\) 贡献一条等差竖线；合并有限多个 \(j\) 得 (1)。
计数上界 (2) 由每条竖线在高度窗 \(|\Im(s)|\le T\) 内最多容纳 \(O(T)\) 个点立得。
\end{proof}

\begin{remark}[为什么这是通往 RH 的结构性门槛]\label{rem:single-scale-dirichlet-obstruction}
黎曼 \(\zeta(s)\) 的非平凡零点满足 \(N_\zeta(T)\sim \frac{T}{2\pi}\log\frac{T}{2\pi}\)，其零点密度随 \(T\) 以 \(\log T\) 增长。
而单尺度 Dirichlet 化的任何有理模型（如 \(\det(I-L^{-s}M)^{-1}\) 及其有限乘积/商）其零点/极点计数至多为 \(O(T)\)（推论 \ref{cor:single-scale-dirichlet-lattice}），因此在“零点密度阶”上就与真实 \(\zeta\) 的现象不兼容。
这解释了为什么引入多尺度（例如下面的连续底混合）是从 toy-RH 走向 RH 研究形态的必要升级。
\end{remark}

\paragraph{单尺度完成化梳齿的计数常数、相关退化与不可公度尺度的去周期化}
据内部文稿 \cite{Internal2025ResolutionFoldingPhiPi} 所固定的“相位坐标门—有限窗审计—orbit–Levi 刚性”并置口径，在不引入额外公理的前提下，单尺度完成化把零点/极点集合压缩为有限条等差竖线簇；该压缩同时给出一类由刚性内生确定的规范归一化机制，并在解析侧表现为梳齿结构的计数常数与相关退化。
另一方面，黄金底 Mellin 指纹（定理 \ref{thm:gm-mellin-meromorphic-fingerprint} 及其后续推论）诱导另一套基本虚周期为 \(2\pi i/\log\varphi\) 的梳齿；除非两尺度可公度，否则两类梳齿不共享任何非零纯虚周期。更一般地，仅需叠加两不可公度底即可摧毁严格虚周期，从而在结构层面解释“从有限晶格到非晶格”所需的最小离散去周期化机制。

\begin{proposition}[单尺度完成化零点计数的主项常数：由 \texorpdfstring{$\deg P_L$}{deg PL} 刚性确定]\label{prop:xi-single-scale-completed-zero-count-leading-constant}
\leanverified{paper\_xi\_single\_scale\_completed\_zero\_count\_leading\_constant}
固定 \(L>1\) 并沿用 \S\ref{subsec:xi-from-self-dual} 的完成化记号：
令
\(
P_L(S):=\widehat\Delta(L^{-1/2},S)\in\RR[S]
\)
且 \(d:=\deg P_L\)，并令
\[
\Phi_L(y):=y^{d}P_L(y+y^{-1})\in\RR[y]
\]
如推论 \ref{cor:xi-rh-unit-circle-reciprocal}。
记 \(N_{\Xi}(T;\sigma_1,\sigma_2)\) 为条带 \(\sigma_1\le \Re(s)\le \sigma_2\) 内、满足 \(|\Im(s)|\le T\) 的 \(\Xi_{\Delta,L}\) 零点个数（按重数计）。
则存在常数 \(C(\sigma_1,\sigma_2)\) 使得
\[
N_{\Xi}(T;\sigma_1,\sigma_2)
=\frac{\log L}{\pi}\,T\cdot
\#\!\left\{
y\in\CC:\ \Phi_L(y)=0,\ 
\sigma_1\le \tfrac12+\frac{\log|y|}{\log L}\le \sigma_2
\right\}
+O(1)\qquad(T\to\infty).
\]
特别地，在全条带 \((\sigma_1,\sigma_2)=(-\infty,\infty)\) 上有
\[
N_{\Xi}(T)
=\frac{2d\,\log L}{\pi}\,T+O(1).
\]
\end{proposition}

\begin{proof}
由 \(\Phi_L(y)=0\) 当且仅当 \(P_L(y+y^{-1})=0\)（且 \(\Phi_L(0)=\mathrm{lc}(P_L)\neq 0\)），并令 \(y:=L^{s-1/2}\)（命题 \ref{prop:xi-rh-interval-criterion} 的证明），得到
\[
\Xi_{\Delta,L}(s)=0\quad\Longleftrightarrow\quad \Phi_L\!\bigl(L^{s-1/2}\bigr)=0.
\]
因此 \(\Phi_L\) 的每个根 \(y\in\CC^\times\)（计重数）诱导一条零点塔
\[
s=\frac12+\frac{\log y+2\pi i k}{\log L},\qquad k\in\ZZ,
\]
其中 \(\log\) 取任一支，\(2\pi i k\) 表示复对数的多值性。
对固定 \(y\)，该等差列在高度窗 \(|\Im(s)|\le T\) 内的点数为 \(\frac{\log L}{\pi}T+O(1)\)；
并且条带约束等价于
\(
\sigma_1\le \frac12+\frac{\log|y|}{\log L}\le \sigma_2
\)。
对所有根求和即得首式。
在全条带上根数为 \(\deg\Phi_L=2d\)，从而得到
\(
N_{\Xi}(T)=\frac{2d\,\log L}{\pi}T+O(1)
\)。
\end{proof}

\begin{proposition}[单尺度完成化的二点相关退化为 Dirac 梳齿：归一化间隔下的纯点谱]\label{prop:xi-single-scale-two-point-correlation-comb}
\leanverified{paper\_xi\_single\_scale\_two\_point\_correlation\_comb}
保持命题 \ref{prop:xi-single-scale-completed-zero-count-leading-constant} 的设定，并令 \(\mathcal Z\) 为 \(\Xi_{\Delta,L}\) 的零点多重集。
记虚步长
\(
\Delta_t(L):=2\pi/\log L
\)。
对任意紧支撑连续函数 \(\psi\in C_c(\RR)\)，定义归一化二点统计
\[
\mathcal R_T(\psi)
:=\frac{1}{N_{\Xi}(T)}
\sum_{\substack{s,s'\in\mathcal Z\\ |\Im(s)|,|\Im(s')|\le T}}
\psi\!\left(\frac{\Im(s)-\Im(s')}{\Delta_t(L)}\right).
\]
则极限 \(\lim_{T\to\infty}\mathcal R_T(\psi)\) 存在，且为有限个平移 Dirac 梳齿的线性组合；
更具体地，存在有限集 \(A\subset[0,1)\) 与非负权重 \((w_a)_{a\in A}\)（由 \(\Phi_L\) 根塔的相位偏移决定），使得
\[
\lim_{T\to\infty}\mathcal R_T(\psi)
=\sum_{a\in A} w_a\sum_{n\in\ZZ}\psi(n+a).
\]
\end{proposition}

\begin{proof}
由命题 \ref{prop:xi-single-scale-completed-zero-count-leading-constant}，\(\mathcal Z\) 是有限多条等差竖线塔（按重数计）的并。
对任一根 \(y\) 写 \(y=|y|e^{i\vartheta}\) 且令 \(\alpha:=\vartheta/(2\pi)\in[0,1)\)，则相应零点塔的虚部集合为
\[
\Im(s)=\Delta_t(L)\,(\ZZ+\alpha).
\]
因此任取两条塔，其虚部差集落在
\(
\Delta_t(L)\,(\ZZ+(\alpha-\beta))
\)
上；归一化后即 \(\ZZ+(\alpha-\beta)\)。
对有限多塔求和得到有限集 \(A\)（由所有相位偏移差的同余类组成）。
又由于每条塔在高度窗内具有均匀线性计数，紧支撑测试函数 \(\psi\) 仅采样有限个整数平移，从而每一对塔贡献的极限权重存在；
对所有塔对加总即得所示形式。
\end{proof}

\begin{proposition}[两类梳齿尺度的不可公度性：除 \texorpdfstring{$\log L/\log\varphi\in\QQ$}{log L/log phi in Q} 外无共同虚周期]\label{prop:xi-comb-incommensurability-logL-logphi}
\leanverified{paper\_xi\_comb\_incommensurability\_logL\_logphi}
设 \(L>1\)，并令单尺度 Dirichlet 化的基本虚周期为
\(
T_L:=2\pi i/\log L
\)
（命题 \ref{prop:single-scale-dirichlet-imag-period}）。
黄金底 Mellin 变量 \(z=\varphi^{-s}\) 的基本虚周期为
\(
T_\varphi:=2\pi i/\log\varphi
\)。
则存在非零 \(T\in i\RR\) 同时属于 \(T_L\ZZ\) 与 \(T_\varphi\ZZ\) 当且仅当 \(\log L/\log\varphi\in\QQ\)，等价地 \(L=\varphi^{p/q}\)（某些互素整数 \(p,q\)）。
特别地，对任意整数 \(L\ge 2\)，有 \(\log L/\log\varphi\notin\QQ\)，从而两类梳齿不可能具有非零共同周期。
\end{proposition}

\begin{proof}
若 \(T\in T_L\ZZ\cap T_\varphi\ZZ\)，则存在非零整数 \(m,n\) 使
\(
T=2\pi i m/\log L=2\pi i n/\log\varphi
\)，故 \(\log L/\log\varphi=m/n\in\QQ\)。
反向由同式直接构造公共周期。
若 \(L\in\ZZ_{\ge2}\) 且 \(\log L/\log\varphi=p/q\)，则 \(L^q=\varphi^p\)。
左端为有理数而右端为二次无理数（\(p\neq 0\) 时 \(\varphi^p\notin\QQ\)），矛盾。
\end{proof}

\begin{theorem}[双底叠加的去周期化：两不可公度尺度足以摧毁严格虚周期]\label{thm:xi-two-base-deperiodization-no-imag-period}
\leanverified{paper\_xi\_two\_base\_deperiodization\_no\_imag\_period}
取 \(L_1,L_2>1\) 且 \(\log L_1/\log L_2\notin\QQ\)。
令 \(g,h\) 为一变量亚纯函数。
进一步假设二者均不具有任何非平凡的单位根乘法不变性：不存在 \(\omega\neq 1\) 与整数 \(n\ge 2\) 使 \(\omega^n=1\) 且 \(g(\omega z)\equiv g(z)\) 或 \(h(\omega z)\equiv h(z)\)。
定义
\[
f(s):=g(L_1^{-s})+h(L_2^{-s}).
\]
则 \(f\) 不存在任何非零纯虚周期 \(T\in i\RR\setminus\{0\}\)。
\end{theorem}

\begin{proof}
反设存在 \(T=i\tau\neq 0\) 使 \(f(s+T)\equiv f(s)\)。
令 \(\omega_1:=e^{-T\log L_1}\)、\(\omega_2:=e^{-T\log L_2}\)，则 \(|\omega_1|=|\omega_2|=1\) 且
\[
g(\omega_1 L_1^{-s})-g(L_1^{-s})=h(L_2^{-s})-h(\omega_2 L_2^{-s})
\]
作为 \(s\) 的亚纯恒等式成立。
\par
取 \(r>0\) 使 \(g,h\) 在穿孔圆盘 \(0<|z|<r\) 上作 Laurent 展开
\[
g(z)=\sum_{m\ge -M_g} a_m z^m,\qquad
h(z)=\sum_{n\ge -M_h} b_n z^n
\qquad(0<|z|<r),
\]
其中 \(M_g,M_h\ge 0\)。
取 \(\sigma_0\) 足够大使得当 \(\Re(s)\ge \sigma_0\) 时 \(0<|L_1^{-s}|<r\) 且 \(0<|L_2^{-s}|<r\)。
代入 \(z=L_i^{-s}\) 得
\[
g(L_1^{-s})=\sum_{m\ge -M_g} a_m L_1^{-ms},\qquad
h(L_2^{-s})=\sum_{n\ge -M_h} b_n L_2^{-ns}
\qquad(\Re(s)\ge \sigma_0).
\]
由周期性 \(f(s+T)\equiv f(s)\) 得
\[
0=f(s+T)-f(s)
=\sum_{m\ge -M_g} a_m(\omega_1^m-1)L_1^{-ms}
\;+\sum_{n\ge -M_h} b_n(\omega_2^n-1)L_2^{-ns}.
\]
由于 \(\log L_1/\log L_2\notin\QQ\)，集合 \(\{m\log L_1:\ m\in\ZZ\}\cup\{n\log L_2:\ n\in\ZZ\}\) 中除 \(0\) 外无重合；
由 Dirichlet 型指数级数（或指数多项式与其绝对收敛尾项）的唯一性，必对所有 \(m,n\) 有
\(
a_m(\omega_1^m-1)=0
\) 与
\(
b_n(\omega_2^n-1)=0
\)。
因此对所有 \(a_m\neq 0\) 的 \(m\) 有 \(\omega_1^m=1\)，对所有 \(b_n\neq 0\) 的 \(n\) 有 \(\omega_2^n=1\)。
若对所有 \(m\neq 0\) 均有 \(a_m=0\)，则 \(g\) 为常数，从而对任意单位根 \(\omega\neq 1\) 有 \(g(\omega z)\equiv g(z)\)，与假设矛盾；
故存在 \(m\neq 0\) 使 \(a_m\neq 0\)，从而 \(\omega_1\) 为单位根；同理 \(\omega_2\) 亦为单位根。
并且在 \(0<|z|<r\) 上有
\[
g(\omega_1 z)
=\sum_{m\ge -M_g} a_m(\omega_1 z)^m
=\sum_{m\ge -M_g} a_m\omega_1^m z^m
=\sum_{m\ge -M_g} a_m z^m
=g(z),
\]
由亚纯延拓得 \(g(\omega_1 z)\equiv g(z)\)；
同理 \(h(\omega_2 z)\equiv h(z)\)。
由单位根不变性假设只能有 \(\omega_1=\omega_2=1\)。
\par
于是 \(T\in (2\pi i/\log L_1)\ZZ\) 且 \(T\in (2\pi i/\log L_2)\ZZ\)。
由 \(\log L_1/\log L_2\notin\QQ\) 得两格交仅含 \(0\)，与 \(T\neq 0\) 矛盾。
\end{proof}


\paragraph{尺度 Mellin 叠加完成化（连续底混合）}
文稿已证明：在自对偶条件下，对任意底 \(L>1\) 都有严格函数方程（命题 \ref{prop:xi-delta-L-functional-equation}）与临界线实值性（命题 \ref{prop:xi-delta-L-critical-line-real}）。
因此可把 \(L\) 视为一个连续尺度参数，并对其做 Mellin 型叠加。

\begin{definition}[尺度叠加完成化]\label{def:xi-scale-mellin-superposition}
取一个实值权函数 \(w:[0,\infty)\to\RR\) 使下述积分对所关心的 \(s\) 收敛（例如 \(w\in L^1(0,\infty)\) 且具有足够衰减）。
定义
\[
\boxed{\ \mathcal X_{\Delta,w}(s):=\int_{1}^{\infty} w(\log L)\,\Xi_{\Delta,L}(s)\,\frac{dL}{L}\ }.
\]
作变量代换 \(x=\log L\) 得等价形式
\[
\mathcal X_{\Delta,w}(s)=\int_0^\infty w(x)\,\Delta\!\bigl(e^{-sx},e^{(2s-1)x}\bigr)\,dx.
\]
\end{definition}

\begin{proposition}[尺度叠加保持函数方程与临界线实值性]\label{prop:xi-scale-mellin-fe-real}
\leanverified{paper\_xi\_scale\_mellin\_fe\_real}
在定义 \ref{def:xi-scale-mellin-superposition} 的条件下，有严格恒等式
\[
\boxed{\ \mathcal X_{\Delta,w}(s)=\mathcal X_{\Delta,w}(1-s)\ }.
\]
并且对任意 \(t\in\RR\) 有 Hardy 型实值化
\[
\boxed{\ \mathcal X_{\Delta,w}\!\left(\tfrac12+it\right)\in\RR\ }.
\]
\end{proposition}
\begin{proof}
对每个 \(L\) 有 \(\Xi_{\Delta,L}(s)=\Xi_{\Delta,L}(1-s)\)，逐点代入积分即得函数方程。
当 \(s=\tfrac12+it\) 时，命题 \ref{prop:xi-delta-L-critical-line-real} 给出对每个 \(L\) 都有 \(\Xi_{\Delta,L}(s)\in\RR\)；又因 \(w\) 为实值，故积分为实数。
\end{proof}

\begin{remark}[动机：连续尺度混合打破有限竖向晶格刚性]\label{rem:xi-scale-mellin-break-lattice}
对固定有限维 \(M\) 的 \(\det(I-\lambda^{-s}M)\)（或有限个此类对象的有限乘积/商），其零极点只能落在有限条等差竖线簇上（推论 \ref{cor:single-scale-dirichlet-lattice}；亦可参见附录 \ref{app:real-input-40-finite-rh} 的谱--极点字典），因此其计数函数至多为 \(O(T)\)。
而 \(\mathcal X_{\Delta,w}\) 是对连续尺度族 \(L\in(1,\infty)\) 的叠加：在临界线 \(s=\tfrac12+it\) 上
\[
\mathcal X_{\Delta,w}\!\left(\tfrac12+it\right)=\int_0^\infty w(x)\,\Delta\!\bigl(e^{-x/2}e^{-itx},e^{2itx}\bigr)\,dx,
\]
这是一个以 \(x\) 为对数尺度、以 \(t\) 为频率的振荡积分；其可调旋钮是 \(w\) 的尾部（决定叠加了多大的频率窗）。
这为“从有限晶格到可能的非晶格零点集”的跃迁提供了一个结构上自然的通道。
\end{remark}

\begin{proposition}[有限频带/有限尺度的零点密度上界：\texorpdfstring{$N(T)=O(T)$}{N(T)=O(T)}]\label{prop:xi-scale-mellin-zero-density-upper}
设 \(\Delta(z,u)\) 在 unitary slice 上具有有限带宽：存在整数 \(N,K\ge 0\) 以及系数 \(a_{j,k}\in\CC\)，使得
\[
\Delta(z,u)=\sum_{j=0}^{N}\sum_{k=-K}^{K} a_{j,k}\,z^j u^k
\qquad(\text{有限和}).
\]
再设权函数 \(w\) 具有紧支集 \(\mathrm{supp}(w)\subseteq [0,X]\)（某个 \(X<\infty\)）。
令
\[
F(t):=\mathcal X_{\Delta,w}\!\left(\tfrac12+it\right)\qquad(t\in\RR),
\]
并令 \(N_F(T)\) 表示区间 \([-T,T]\) 上 \(F(t)\) 的实零点数（按重数计）。
则：
\begin{enumerate}
  \item \(F\) 作为 \(t\)-变量的函数可延拓为一个指数类型有界的整函数（其类型由 \(X\) 与带宽 \(N,K\) 控制）。
  \item 因而零点计数满足线性上界
  \[
  \boxed{\ N_F(T)=O(T)\qquad(T\to\infty)\ }.
  \]
\end{enumerate}
\end{proposition}
\begin{proof}
在假设下，取 \(s=\tfrac12+it\) 并写 \(u=e^{2itx}\)、\(z=e^{-x/2}e^{-itx}\)，则
\[
F(t)=\int_0^X w(x)\sum_{j=0}^{N}\sum_{k=-K}^{K} a_{j,k}\,e^{-jx/2}\,e^{i(2k-j)tx}\,dx
=\sum_{j,k}a_{j,k}\,\widehat g_{j}( (2k-j)t),
\]
其中 \(g_j(x):=w(x)e^{-jx/2}\,\mathbf{1}_{[0,X]}(x)\in L^1(\RR)\)，而 \(\widehat g_j(\xi):=\int_{\RR} g_j(x)e^{i\xi x}\,dx\) 为其 Fourier 变换。
由 \(\mathrm{supp}(g_j)\subseteq[0,X]\)，Paley--Wiener 型结论给出 \(\widehat g_j\) 为指数类型 \(\le X\) 的整函数，从而 \(F\) 亦为指数类型有限的整函数（类型可被 \(X\max_{j,k}|2k-j|\le X(2K+N)\) 控制）。
指数类型整函数的零点分布满足线性密度上界，故 \(N_F(T)=O(T)\)。
\end{proof}

\begin{proposition}[连续尺度叠加的竖向非周期性：除零周期外不存在共同纯虚周期]\label{prop:xi-scale-mellin-no-nonzero-imag-period}
设 \(\Delta\) 为有限项和
\(
\Delta(z,u)=\sum_{j,k}a_{j,k}z^j u^k
\)，并取实值权函数 \(w\) 满足 \(\mathrm{supp}(w)\subseteq[0,X]\)。
定义尺度叠加完成化
\(
\mathcal X_{\Delta,w}(s)=\int_0^\infty w(x)\Delta(e^{-sx},e^{(2s-1)x})\,dx
\)
如定义 \ref{def:xi-scale-mellin-superposition}。
对每个整数 \(m\) 定义
\[
b_m(x):=\sum_{j-2k=m}a_{j,k}e^{-kx},
\]
并假设存在某个 \(m\neq 0\) 使得 \(w(x)b_m(x)\) 在某个区间上非零于正测度集，且 \(\mathcal X_{\Delta,w}\not\equiv 0\)。
则 \(\mathcal X_{\Delta,w}\) 不存在任何非零纯虚周期：若对某个 \(T\in i\RR\) 有
\(
\mathcal X_{\Delta,w}(s+T)\equiv \mathcal X_{\Delta,w}(s)
\)，则必有 \(T=0\)。
\leanverified{paper\_xi\_scale\_mellin\_no\_nonzero\_imag\_period}
\end{proposition}

\begin{proof}
反设存在 \(T=i\tau\neq 0\)（\(\tau\in\RR\)）使 \(\mathcal X_{\Delta,w}(s+i\tau)\equiv \mathcal X_{\Delta,w}(s)\)。
固定任意 \(\sigma\in\RR\)，并令
\[
F_\sigma(t):=\mathcal X_{\Delta,w}(\sigma+\mathrm{i}t)\qquad(t\in\RR).
\]
由 \(\mathrm{supp}(w)\subseteq[0,X]\) 与 \(\Delta\) 为有限项和，展开得
\[
F_\sigma(t)
=\sum_{j,k}a_{j,k}\int_0^X w(x)e^{-kx}e^{-\sigma(j-2k)x}\,e^{-\mathrm{i}t(j-2k)x}\,dx,
\]
故 \(F_\sigma\) 为若干 \(L^1\) 函数的 Fourier 变换之有限线性组合。
由 Riemann--Lebesgue 引理，必有 \(F_\sigma(t)\to 0\) 当 \(|t|\to\infty\)。
\par
另一方面，周期性给出对所有 \(t\) 有 \(F_\sigma(t+\tau)=F_\sigma(t)\)。
令 \(n\to\infty\)，则
\[
F_\sigma(t)=F_\sigma(t+n\tau)\longrightarrow 0,
\]
从而 \(F_\sigma(t)\equiv 0\)。
由于 \(\sigma\) 任意，得 \(\mathcal X_{\Delta,w}\) 在整个竖直直线上恒为零；由解析函数唯一性定理推出 \(\mathcal X_{\Delta,w}\equiv 0\)，与假设矛盾。
故 \(\tau=0\)，即 \(T=0\)。
\end{proof}

\begin{proposition}[有限频带与紧支撑下的 Cartwright 指数型：显式类型上界与线性零点密度]\label{prop:xi-scale-mellin-cartwright-explicit-type}
在命题 \ref{prop:xi-scale-mellin-zero-density-upper} 的设定下，进一步令
\[
B:=\max\{|j-2k|:\ a_{j,k}\neq 0\},
\qquad
\mathrm{supp}(w)\subseteq[0,X].
\]
记
\(
F(t):=\mathcal X_{\Delta,w}\!\left(\tfrac12+it\right)
\)
与其零点计数 \(N_F(T)\) 如命题 \ref{prop:xi-scale-mellin-zero-density-upper}。
则 \(F\) 属 Cartwright 类整函数，且其指数类型不超过 \(XB\)；
从而临界线零点计数满足显式上界
\[
N_F(T)\le \frac{2XB}{\pi}\,T+o(T)\qquad(T\to\infty).
\]
\end{proposition}

\begin{proof}
命题 \ref{prop:xi-scale-mellin-zero-density-upper} 的证明已给出分解
\[
F(t)=\sum_{j,k}a_{j,k}\,\widehat g_j\bigl((2k-j)t\bigr),
\qquad
g_j(x):=w(x)e^{-jx/2}\mathbf 1_{[0,X]}(x).
\]
由于 \(\mathrm{supp}(g_j)\subseteq[0,X]\)，故 \(\widehat g_j\) 为指数类型 \(\le X\) 的整函数；复合缩放 \(t\mapsto (2k-j)t\) 使指数类型至多乘以 \(|2k-j|\)，因此 \(F\) 的指数类型不超过 \(X\max_{j,k}|2k-j|=XB\)。
又因 \(\widehat g_j\) 在实轴上由 \(\|g_j\|_{L^1}\) 统一界定，\(F\) 在实轴上亦有界；
故 \(F\) 满足 Cartwright 类的定义性增长条件。
对 Cartwright 类整函数应用 Jensen--Cartwright 的零点密度估计，即得
\(
N_F(T)\le \frac{2\,\mathrm{type}(F)}{\pi}T+o(T)\le \frac{2XB}{\pi}T+o(T)
\)。
\end{proof}

\begin{corollary}[\texorpdfstring{$T\log T$}{T log T} 零点密度的必要尺度窗：紧支撑叠加无法达到]\label{cor:xi-scale-mellin-window-necessary-tlogt}
\leanverified{paper\_xi\_scale\_mellin\_window\_necessary\_tlogt}
在命题 \ref{prop:xi-scale-mellin-cartwright-explicit-type} 的有限带宽设定下，考虑一族权重 \(w_T\) 满足
\(
\mathrm{supp}(w_T)\subseteq[0,X(T)]
\)
且带宽参数 \(B\) 保持不变。
令
\(
F_T(t):=\mathcal X_{\Delta,w_T}(\tfrac12+it)
\)
且 \(N_{F_T}(T)\) 为其在 \([-T,T]\) 上的零点数（按重数计）。
则有
\[
N_{F_T}(T)\le \frac{2B}{\pi}\,X(T)\,T+o\!\bigl(X(T)T\bigr)\qquad(T\to\infty).
\]
因此若 \(X(T)=o(\log T)\)，则必有 \(N_{F_T}(T)=o(T\log T)\)。
\end{corollary}

\begin{proof}
对每个 \(T\) 将命题 \ref{prop:xi-scale-mellin-cartwright-explicit-type} 应用于 \(w=w_T\)（此时 \(X\) 替换为 \(X(T)\)）即得首式。
若 \(X(T)=o(\log T)\)，则 \(\frac{2B}{\pi}X(T)T=o(T\log T)\)，从而 \(N_{F_T}(T)=o(T\log T)\)。
\end{proof}

\begin{corollary}[\texorpdfstring{$T\log T$}{T log T} 零点密度的尺度宽度下界：必须混合到 \texorpdfstring{$X(T)\gtrsim \frac{1}{2B}\log T$}{X(T) gtrsim (1/2B) log T}]\label{cor:xi-scale-mellin-window-width-lower-bound-1over2B}
沿用推论 \ref{cor:xi-scale-mellin-window-necessary-tlogt} 的设定，并记 \(\widetilde N_{F_T}(T):=N_{F_T}(T)\)。
若目标是使对称零点密度达到黎曼型量级
\[
\widetilde N_{F_T}(T)\ \asymp\ \frac{T}{\pi}\log T,
\]
则必有
\[
\liminf_{T\to\infty}\frac{X(T)}{\log T}\ \ge\ \frac{1}{2B}.
\]
\end{corollary}

\begin{proof}
由推论 \ref{cor:xi-scale-mellin-window-necessary-tlogt} 得
\(
\widetilde N_{F_T}(T)\le \frac{2B}{\pi}X(T)T+o(X(T)T)
\)。
与 \(\widetilde N_{F_T}(T)\asymp (T/\pi)\log T\) 对比，得到 \(X(T)\gtrsim (1/(2B))\log T\)，即所示 \(\liminf\) 下界。
\end{proof}

\begin{corollary}[底参数的幂级范围门槛：逼近零点密度强制 \texorpdfstring{$L\lesssim T^{1/(2B)}$}{L lesssim T^{1/(2B)}} 的混合范围]\label{cor:xi-scale-mellin-base-range-threshold}
在推论 \ref{cor:xi-scale-mellin-window-width-lower-bound-1over2B} 的设定下，令 \(L=e^{x}\)。
权函数支撑 \(x\in[0,X(T)]\) 等价于底参数混合范围
\(
1\le L\le e^{X(T)}
\)。
若 \(\widetilde N_{F_T}(T)\asymp (T/\pi)\log T\)，则必有
\[
e^{X(T)}\ \ge\ T^{\,\frac{1}{2B}+o(1)}.
\]
特别地，在最小带宽 \(B=1\) 的情形，若不混合到 \(L\ge T^{1/2+o(1)}\) 的范围，则在零点密度阶上不可能达到黎曼型量级。
\end{corollary}

\begin{proof}
由推论 \ref{cor:xi-scale-mellin-window-width-lower-bound-1over2B} 得 \(X(T)\ge (\frac{1}{2B}+o(1))\log T\)。
两端取指数即得 \(e^{X(T)}\ge T^{1/(2B)+o(1)}\)。
\end{proof}


\begin{remark}[密度门槛：逼近 \texorpdfstring{$T\log T$}{T log T} 必须打破“有限频带”]\label{rem:xi-scale-mellin-density-threshold}
命题 \ref{prop:xi-scale-mellin-zero-density-upper} 表明：若 \(\Delta\) 的 \(u\)-带宽固定且 \(w\) 在对数尺度上具有有效截断（紧支集或同阶意义下的硬截断），则即便引入连续尺度混合，临界线零点计数仍至多为 \(O(T)\)。
因此，若目标是逼近 \(\zeta\) 的零点计数增长阶（\(N_\zeta(T)\asymp T\log T\)），则必须至少引入一条硬升级机制：
\begin{enumerate}
  \item \textbf{无界尺度尾部}：使得对频率 \(T\) 的“有效尺度上界” \(X(T)\) 随 \(T\) 增长（例如 \(X(T)\asymp \log T\) 的尺度律）；或
  \item \textbf{随尺度增长的复杂度}：让 \(\Delta\) 的带宽/状态数（例如 \(N,K\)）随尺度参数增长，而非固定为有限常数。
\end{enumerate}
该门槛把“打破竖向晶格”进一步细化为“打破有限频带”：仅靠有限尺度混合并不足以在密度阶上模拟 \(\zeta\)-型零点集合。
\end{remark}

\begin{remark}[\texorpdfstring{$\log T$}{log T} 预算：尺度尾窗宽 \texorpdfstring{$X(T)$}{X(T)} 与相位模态数 \texorpdfstring{$B(T)$}{B(T)} 的同阶门槛]\label{rem:xi-scale-mellin-logT-budget}
命题 \ref{prop:xi-scale-mellin-zero-density-upper} 的证明实际上给出一条更定量的结果：对任意 \(X<\infty\)，若把尺度积分硬截断为
\[
\mathcal X_{\Delta,w}^{(\le X)}(s):=\int_{0}^{X} w(x)\,\Delta\!\bigl(e^{-sx},e^{(2s-1)x}\bigr)\,dx,
\]
则 \(t\mapsto \mathcal X_{\Delta,w}^{(\le X)}(\tfrac12+it)\) 的指数类型至多为 \(O_\Delta(X)\)，从而其高度窗零点数满足
\[
\boxed{\;N^{(\le X)}(T)=O_\Delta(X\,T)\; }.
\]
因此，若对给定高度 \(T\) 而言，\(\mathcal X_{\Delta,w}(\tfrac12+it)\) 的“有效尺度窗宽”必须增长到 \(X=X(T)\)（例如：忽略 \(x>X(T)\) 的尾项会在该高度窗内造成不可接受的误差），则想要达到黎曼型密度 \(N(T)\asymp T\log T\) 就强迫
\[
\boxed{\;X(T)\asymp \log T\; }.
\]
\par
同一 \(\log T\) 预算也出现在“相位/角色轴”的最小无限机制中：在有限维临界线梳齿框架内，若试图通过引入随尺度增长的临界圆模态数 \(B=B(T)\)（例如通过循环 lift 的 Artin 因子化把单位根相位注入并令 \(q=q(T)\to\infty\)）来实现 \(N(T)\asymp T\log T\)，则必要地需有 \(B(T)\asymp (\log T)/(\log\lambda)\)（推论 \ref{cor:critical-circle-mode-growth}）。
因此在本文框架下，任何“逼近 \(\zeta\) 型零点密度”的最小无限来源，本质上都需要在某个方向注入 \(\Theta(\log T)\) 的自由度：要么在尺度尾部释放 \(\log T\) 级的窗宽，要么在相位/角色方向释放 \(\log T\) 级的模态数；命题 \ref{prop:xi-scale-mellin-bessel-kernel} 给出的 Bessel--cosh 核提供了一个显式可校准的实现样例。
\end{remark}

\begin{proposition}[Bessel--cosh 尺度核给出 \texorpdfstring{$T\log T$}{T log T} 临界线零点密度与 \texorpdfstring{$2\pi$}{2pi} 常数锚点]\label{prop:xi-scale-mellin-bessel-kernel}
取参数 \(x_0>0\)，并定义偶权核
\[
w_{x_0}(t):=\exp\!\bigl(-x_0\cosh t\bigr)\qquad(t\ge 0).
\]
则 Paris 的积分表示给出（\(\nu\in\RR\)）
\[
\boxed{\;
K_{i\nu}(x_0)=\int_{0}^{\infty}e^{-x_0\cosh t}\cos(\nu t)\,dt
\;} \qquad\text{\cite{Paris2022KiNuZeros}},
\]
从而沿临界线的“Bessel 完成模块”
\[
\Xi_{\mathrm{Bessel},x_0}(s):=K_{\frac{s-\frac12}{2}}(x_0)
\]
满足严格对称性与 Hardy 实值性
\[
\boxed{\ \Xi_{\mathrm{Bessel},x_0}(s)=\Xi_{\mathrm{Bessel},x_0}(1-s)\ },\qquad
\boxed{\ \Xi_{\mathrm{Bessel},x_0}\!\left(\tfrac12+it\right)\in\RR\ }.
\]
并且其临界线零点由 \(K_{i\nu}(x_0)\) 的 \(\nu\)-零点给出：若 \(K_{i\nu}(x_0)=0\) 且 \(\nu>0\)，则
\(\Xi_{\mathrm{Bessel},x_0}(\tfrac12+2i\nu)=0\)。
记
\[
N_{\mathrm{Bessel},x_0}(T):=\#\Bigl\{\,0<\gamma\le T:\ \Xi_{\mathrm{Bessel},x_0}\!\left(\tfrac12+i\gamma\right)=0\,\Bigr\},
\]
则 Paris 的大阶渐近方程 \cite{Paris2022KiNuZeros} 推出
\[
\boxed{\;
N_{\mathrm{Bessel},x_0}(T)
=\frac{T}{2\pi}\log\!\Bigl(\frac{T}{e\,x_0}\Bigr)+O(1)\qquad(T\to\infty).\;}
\]
特别地，当取 \(\boxed{x_0=2\pi}\) 时，
\[
\boxed{\;
N_{\mathrm{Bessel},2\pi}(T)
=\frac{T}{2\pi}\log\!\Bigl(\frac{T}{2\pi}\Bigr)-\frac{T}{2\pi}+O(1),\;}
\]
其前两项常数与黎曼--冯\,曼戈尔特零点计数主项完全同型。
\end{proposition}
\begin{proof}
积分表示与 \(K_\nu(x)=K_{-\nu}(x)\) 给出对称性；当 \(s=\tfrac12+it\) 时阶数为纯虚 \(i t/2\)，故 \(K_{i(t/2)}(x_0)\in\RR\)（\(x_0>0\)），得 Hardy 实值性。
令 \(K_{i\nu}(x_0)\) 的正零点按大小记为 \(\nu_n(x_0)\)；Paris 的零点渐近式给出计数
\(\#\{\nu_n(x_0)\le V\}=\frac{V}{\pi}\log\!\bigl(\frac{2V}{e x_0}\bigr)+O(1)\)。
取 \(V=T/2\) 即得 \(N_{\mathrm{Bessel},x_0}(T)=\#\{\nu_n(x_0)\le T/2\}\) 的结论；令 \(x_0=2\pi\) 立得常数对齐。
\end{proof}
\endinput
\begin{corollary}[混合时间 \texorpdfstring{$\tau_{\mathrm{mix}}(m)\sim\mathrm{const}\cdot m^2$}{tau_mix(m)~const*m^2} 的常数闭式]\label{cor:near1-pole-taumix-m2}
令 \(\rho_m:=\rho(M_{it_m})\) 且定义谱比混合时间
\[
\tau_{\mathrm{mix}}(m):=\frac{1}{-\log(\rho_m/\lambda)}.
\]
则由 \(\Re(P_2(it)-P_2(0))=-\sigma_2^2 t^2/2+O(t^4)\) 得
\[
\boxed{\;
\tau_{\mathrm{mix}}(m)
=\frac{m^2}{2\pi^2\sigma_2^2}\bigl(1+O(m^{-2})\bigr)\; }.
\]
\end{corollary}

\begin{corollary}[Abelian tower 的否定性结果：统一平方根谱隙假设不成立]\label{cor:abelian-tower-no-go}
\leanverified{paper\_abelian\_tower\_no\_go}
设一族同余层 \(G_m\) 含有 \(\ZZ/m\ZZ\) 的阿贝尔分量，使得存在一维单位根扭曲 \(u=\omega_m\to 1\)
对应于某个整数值 cocycle \(F\) 且 \(\sigma_F^2=\mathrm{Var}_\infty(F)>0\)。
则最坏非平凡谱半径比 \(\eta_m\) 满足
\[
\eta_m\ \ge\ \exp\!\Bigl(-\frac{\sigma_F^2}{2}(2\pi/m)^2+O(m^{-4})\Bigr)\ \xrightarrow[m\to\infty]{}\ 1,
\]
从而不存在任何 \(\varepsilon>0\) 使得 \(\eta_m\le \lambda^{-1/2-\varepsilon}\) 对所有 \(m\) 成立。
因此，定理 \ref{thm:auditable-profinite-chebotarev-finite} 中用于推出“平方根级尾项”的统一谱隙假设，
在含有此类阿贝尔模数塔的 profinite 轴上与 near-\(1\) 扭曲慢模的扩散下界不相容。
\end{corollary}

\begin{table}[H]
\centering
\scriptsize
\setlength{\tabcolsep}{6pt}
\caption{Near-1 pole diffusive prediction for the collision twist ($\varphi_2=\mathbf{1}_{\{d=2\}}$) of the real-input 40-state kernel. Constants are computed from $\mu_2=P_2'(0)$, $\sigma_2^2=P_2''(0)$, and $\lambda=\varphi^2$.}
\label{tab:real_input_40_collision_near1_pole_predictions}
\begin{tabular}{r r r}
\toprule
$m$ & pred. $\Re(s_m)$ & pred. $\Im(s_m)$\\
\midrule
5 & 0.944762 & 0.099747\\
7 & 0.971817 & 0.071248\\
10 & 0.986190 & 0.049873\\
20 & 0.996548 & 0.024937\\
50 & 0.999448 & 0.009975\\
\midrule
$c$ & \multicolumn{2}{l}{1.380957}\\
$d$ & \multicolumn{2}{l}{0.498733}\\
$\mu_2$ & \multicolumn{2}{l}{0.0763932023}\\
$\sigma_2^2$ & \multicolumn{2}{l}{0.0673312629}\\
$\log\lambda$ & \multicolumn{2}{l}{0.9624236501}\\
$\frac{\sigma_2^2\log\lambda}{2\mu_2^2}$ & \multicolumn{2}{l}{5.551925}\\
\bottomrule
\end{tabular}
\end{table}


\begin{conjecture}[尺度化核版 GRH/Chebotarev（双极限扩散口径）]\label{conj:scaled-kernel-grh-diffusive}
存在常数 \(c>0\) 与 \(\kappa>0\)，使得对所有模数 \(m\ge 2\) 与所有长度 \(n\ge 1\)，任意由模 \(m\) 寄存器定义的同余 cylinder 事件的等分布误差满足
\[
\boxed{\ \mathrm{Err}(n,m)\ \le\ C(m)\,\exp\!\Bigl(-c\,\frac{n}{m^2}\Bigr)\ },
\]
其中 \(C(m)\) 至多多项式增长，并且主指数 \(c\) 与二次律常数 \(\kappa\) 同阶（在可解析的核族中可由压力曲率/方差密度给出）。
因此在区域 \(m\le n^{1/2-\varepsilon}\)（任意固定 \(\varepsilon>0\)）内得到一致的“平方根级可见主项”。
\end{conjecture}
